\begin{theorem}[the candidate conductor mechanism, with all hypotheses]
\label{thm:candidate-conductor}
\Verified\ Assume:
\begin{enumerate}[label=(H\arabic*)]
\item $X$ is a connected compact smooth complex threefold and
$f:X\to B$ is a nonconstant holomorphic map to a smooth compact curve;
\item $W=f^{-1}(p_0)$ is the full scheme-theoretic fibre, is reduced, and
has the normal-crossings local models (5.3);
\item the normalization $\eta:\widetilde W\to W$ is the degree-six del
Pezzo surface $\mathrm{dP}_6$;
\item $H^2(X;\mathbb Z)=0$.
\end{enumerate}
Then, for every $L\in\operatorname{Pic}(X)$,
\[
H^2(W,(T_X\otimes L)|_W)\ne0
\quad\text{and}\quad
R^2f_*(T_X\otimes L)\ne0.
\tag{5.18}
\]
\end{theorem}

\begin{proof}
Fix $L$ and put
\[
\mathcal F_L=T_X\otimes L,\qquad
A_L=(L^\vee\otimes\omega_X)|_W.
\]
Hypothesis (H4) gives $c_1(L^\vee\otimes\omega_X)=0$.  Pullback therefore
gives $c_1(\eta^*A_L)=0$.  The surface $\mathrm{dP}_6$ is the blow-up of
$\mathbb P^2$ at three noncollinear points; the standard blow-up calculation
gives $H^1(\mathrm{dP}_6,\mathcal O)=0$
\cite[Chapter~V]{Hartshorne1977}.  In the exponential sequence, the kernel
of
\[
c_1:\operatorname{Pic}(\mathrm{dP}_6)\longrightarrow
H^2(\mathrm{dP}_6;\mathbb Z)
\]
is the image of $H^1(\mathrm{dP}_6,\mathcal O)$ and is therefore zero.
Consequently
\[
\eta^*A_L\simeq\mathcal O_{\widetilde W}.
\tag{5.19}
\]
\Cref{prop:conductor} now gives
\[
H^0(W,\Omega^1_X|_W\otimes A_L)\ne0.
\tag{5.20}
\]

Because $W=f^*(p_0)$, one has
$\mathcal O_X(W)=f^*\mathcal O_B(p_0)$, whose restriction to the fibre is
trivial.  Adjunction therefore gives $\omega_W\simeq\omega_X|_W$.
Apply \cref{thm:cartier-duality} to $E=\mathcal F_L$:
\[
\begin{split}
H^2(W,\mathcal F_L|_W)^\vee
&\simeq H^0(W,(\mathcal F_L|_W)^\vee\otimes\omega_W)\\
&=H^0(W,\Omega^1_X|_W\otimes A_L).
\end{split}
\]
The right side is nonzero by (5.20), proving the first assertion of (5.18).
The same duality theorem gives
$H^3(W,\mathcal F_L|_W)=0$.

Since $X$ is compact, $f$ is proper.  By
\cref{lem:curve-flatness}, $f$ is flat and the vector bundle
$\mathcal F_L$ is $f$-flat.  Thus every hypothesis of
\cref{prop:analytic-base-change} is satisfied at $p_0$, and
\[
(R^2f_*\mathcal F_L)_{p_0}\otimes\mathbb C(p_0)
\simeq H^2(W,\mathcal F_L|_W)\ne0.
\]
The stalk, and hence the coherent direct-image sheaf, is nonzero.  This is
the second assertion of (5.18).
\end{proof}

